\section{Whole image ideals and actual module factorizations}

\begin{proposition}[The whole image ideal is the least factorization ideal]
\label{prop:actual-quasicoherent-image-factorization}
Let $M$ be a quasicoherent module on an arbitrary scheme $X$, let
$\alpha:M\to\mathcal O_X$, and let $I$ be an ideal sheaf. Then
$\operatorname{im}(\alpha)\subseteq I$ if and only if there is a whole
module morphism $b:M\to I$ whose composite with the original inclusion
$I\hookrightarrow\mathcal O_X$ is $\alpha$. Such a morphism is unique.
The forward implication gives a canonical factor, retaining the value
of every original section on every original open.
\lean{EnvelopeProperty.QuasicoherentImageFactorization.factorOfImageLE, EnvelopeProperty.QuasicoherentImageFactorization.factorOfImageLE_inclusion, EnvelopeProperty.QuasicoherentImageFactorization.imageIdeal_le_iff_factor, EnvelopeProperty.QuasicoherentImageFactorization.existsUnique_factor}
\uses{def:fc-quasicoherent-image-module-factor}
\leanok
\end{proposition}
\begin{proof}
Suppose the image ideal is contained in $I$. For a section $m$ on an
arbitrary open $W$, restrict to any affine open $U\subseteq W$.
Naturality identifies the restriction of $\alpha(m)$ with the image
under $\alpha$ of the restriction of $m$ to $U$. This belongs to the
image ideal on $U$, and hence to $I(U)$. Thus $\alpha(m)$ is a section
of the whole original ideal module on $W$. Corestricting the original
linear section maps gives $b$; their naturality follows from that of
$\alpha$, and the composite is the original map on every open.
Conversely, such a factorization sends every section on an affine open
into the same original affine ideal. Every member of the image ideal
has such a preimage, proving containment. The ideal inclusion is
monomorphic, so cancellation proves uniqueness. No affine, finite
generation, Noetherian, or nonempty hypothesis on $X$ is needed.
\end{proof}
